% ATTENTION : ce fichier n'est plus utilisé pour générer page_de_garde.pdf.
% Depuis le 29/09/2026, page_de_garde.pdf est le PDF officiel envoyé par
% l'école (voir _backup_page_de_garde/ pour la version maison précédente,
% conservée ici uniquement à titre d'archive). NE PAS recompiler ce fichier
% isolément : cela écraserait page_de_garde.pdf avec l'ancienne maquette.
\documentclass[12pt,a4paper]{report}
\usepackage[utf8]{inputenc}
\usepackage[T1]{fontenc}
\usepackage[french]{babel}
\usepackage{graphicx}
\usepackage{geometry}
\usepackage{array}
\usepackage{tikz}

% Configuration des marges
\geometry{top=2cm, bottom=2cm, left=2.5cm, right=2.5cm}

\begin{document}

\begin{titlepage}
    \centering
    
    % --- DOUBLE CADRE DE BORDURE EN TIKZ ---
    \begin{tikzpicture}[overlay,remember picture]
        % Cadre extérieur épais
        \draw[line width=3pt]
            ([shift={(0.8cm,-0.8cm)}]current page.north west)
            rectangle
            ([shift={(-0.8cm,0.8cm)}]current page.south east);
        % Cadre intérieur fin
        \draw[line width=1pt]
            ([shift={(0.95cm,-0.95cm)}]current page.north west)
            rectangle
            ([shift={(-0.95cm,0.95cm)}]current page.south east);
    \end{tikzpicture}

    % --- EN-TÊTE ---
    {\bfseries\large REPUBLIQUE DU SENEGAL}\\[0.3cm]
    
    % Emplacement du drapeau (Remplacez par votre fichier image si disponible)
    \begin{center}
        \includegraphics[height=1.8cm]{drapeau_de_senegal.png} % À remplacer
    \end{center}
    \vspace{-0.2cm}
    
    {\small\itshape Un peuple - Un but - Une foi}\\[0.4cm]
    {\bfseries\small Ministère de l'Enseignement Supérieur et de la Recherche}\\[0.1cm]
    {\bfseries\small Direction Générale de l'Enseignement Supérieur Privé}\\[0.1cm]
    {\bfseries\large Université Dakar Bourguiba}\\[0.5cm]
    {\bfseries\large UFR Sciences et Technologies}\\[0.5cm]
    {\bfseries\large Département Génie Informatique}\\[0.5cm]

    
    % Emplacement du logo ISI
    \begin{center}
        \includegraphics[height=2.2cm]{logo_udb.png} % À remplacer
    \end{center}
    \vspace{0.3cm}

    % --- TITRE DU DIPLÔME ---
    {\bfseries\large MEMOIRE}\\[0.2cm]
    {\small Présenté et soutenu par :}\\[0.1cm]
    {\bfseries\large M. BRAHIMA MAMADOU YARANAGORE}\\[0.2cm]
    {\small Pour l'obtention du diplôme de :}\\[0.1cm]
    {\bfseries\large Master Professionnel : Génie Logiciel}\\[0.1cm]

    % --- SECTION SUJET (CADRE STYLE PARCHEMIN/BOÎTE ARRONDIE) ---
    \raggedright{\bfseries\itshape\large Sujet :}\\[0.2cm]
    \centering
    \begin{tikzpicture}
        \node[draw, line width=1pt, rounded corners=5pt, inner sep=15pt, text width=14cm, align=center] {
            {\bfseries\large « Architecture hybride pour la digitalisation des PME d'import-export en Afrique subsaharienne : conception, déploiement et validation sur le cas SOYAFF (corridor Dakar–Bamako) »}
        };
    \end{tikzpicture}
    \vfill

    % --- DATE DE SOUTENANCE ---
    \hfill{\bfseries\small Soutenu à Dakar, le 3 octobre 2026}\\[0.8cm]

    % --- TABLEAU DES MEMBRES DU JURY ---
    {\bfseries\large Membres du Jury}\\[0.2cm]
    \small
    \begin{tabular}{|l|m{6.5cm}|l|}
        \hline
        \textbf{Statut} & \textbf{NOM et Prénom} & \textbf{Grade} \\ \hline
        Président & & \\ \hline
        Superviseur de mémoire & & \\ \hline
        Superviseur de mémoire & & \\ \hline
        Directeur de mémoire & M. Massamba LO & Ingénieur \\ \hline
        Examinateur 1 : & & \\ \hline
        Examinateur 2 : & & \\ \hline
    \end{tabular}
    \vfill

    % --- ANNÉE ACADÉMIQUE ---
    {\bfseries\large Année académique : 2025-2026}

\end{titlepage}

\end{document}
